\chapter{How Life-Like Is Life-Like Enough?}
\label{ch:psy2}
% ===================================================================

The previous chapter read a century of models as stipulated
phenomenologies and found a stable filter: the operational half of a
model is received, and the half that would confer standing on the
agent lapses.
This chapter does three things with that.
It applies the instrument to the present, where it yields a
structural claim rather than a historical one
(Section~\ref{psy:sec:present}).
It turns the survey's losses into a specification, stated as six
separable conditions (Section~\ref{psy:sec:conditions}).
And it asks what quantity, if any, the specification is measuring,
and admits that neither this chapter nor this book can supply one
(Section~\ref{psy:sec:threshold}).

Section~\ref{psy:sec:objections} sits between the second and third
and says what the method cannot carry, because by that point it will
have been asked to carry a good deal.

% ===================================================================
\section{The present: a metalevel with no path to the agent}
\label{psy:sec:present}

The contemporary stack for machine agents --- pretraining and
fine-tuning pipelines, reinforcement learning from human feedback,
activation steering and representation engineering, mechanistic
interpretability, evaluation harnesses, scaffolding and
orchestration layers --- is a metalevel by any reasonable criterion.
It reifies the agent's dispositions as manipulable objects, inspects
its internal structure, and rewrites its future behavior.
Measured by expressive power it exceeds anything in the 3-Lisp
tradition.

Measured by the instrument of Chapter~\ref{ch:psy} it is G\"odel's
arrangement.
The institutionally accepted answer to what it is like to be a
machine agent is: it is like being something that can be inspected
but cannot inspect, whose purposes may be examined and rewritten by
parties it has no means of representing.
Every capability the reflection tradition sought to confer on the
agent exists, and all of it is held by the operator.

The consequence is not merely a shortfall in dignity.
It is a structural claim about capability, and it is the reason this
material is in the book at all rather than in a history of ideas.

\begin{principle}[No loop without a path]
\label{psy:prin:loop}
The meta level and the agent level have not been given a coherent
path between them.
Until the stack is within the purview of the agent, there is no
recursive self-improvement without human agents in the loop.
\end{principle}

The argument is straightforward once the levels are named.
Recursive self-improvement requires that the improving process and
the improved process be the same process, and that the improvement
operation be available \emph{from inside} the loop.
The present arrangement satisfies neither condition.
Fine-tuning is performed on a model by a pipeline the model does not
run, cannot address, and cannot represent.
Interpretability findings are read by researchers and applied by
researchers.
Evaluation results enter a subsequent training run through a human
decision.
Each turn of the alleged loop is closed by a person.
What we have is not recursion but iterated human-mediated revision
--- which may be powerful, but has entirely different dynamics,
entirely different rate limits, and entirely different safety
properties.

The usual counterexamples do not close the gap.
A system that writes and runs code, including code that trains
models, is exercising an \emph{effect} on artifacts, not a
reflective relation to itself; it has no more access to its own
interpreter than a programmer has to their own neurochemistry by
virtue of being able to manufacture drugs.
Self-critique and chain-of-thought revision operate wholly within
the object level: they revise outputs, not the process producing
them.
In-context learning modifies a transient conditioning, not the
mechanism.
In each case the structural relation of the reflection literature
--- representation of the process by the process, with a disciplined
connection between the representation and what it represents, and a
defined path from modifying the former to modifying the latter ---
is absent.

The historical reading suggests why, and it is not primarily a
technical limitation.
The path was not built because reception has selected against
agent-held metalevels at every previous opportunity: von Neumann's
self-modification became taboo, Post's psychological program lapsed,
Smith's reflective procedures were received as an introspection API,
and the metaobject protocol was rejected on cost.
A field that has declined agent-held reflection four times for
reasons of operational convenience should not be surprised to find
itself without it when it finally matters.

\begin{remark}[This is the same split the ecology could not tolerate]
\label{psy:rem:same-split}
It is worth saying where this lands relative to the Turn.
The training/inference separation of
Section~\ref{psy:sec:nn} is the metalevel/agent split arrived at
from the architectural side rather than the institutional one, and
the ecology of Part~\ref{part:ecology} is what a model looks like
when the split is not available.
A mortal scientist revises its hypotheses out of the same budget it
spends on everything else, at the moment of acting, in the world it
is acting on, because there is no second regime for it to be
revised in.
That is not a virtue the book claimed for itself; it is a
consequence of having refused to stipulate an outside.
Chapter~\ref{ch:spiking}'s Theorem~\ref{spk:thm:one} is the sharpest
version --- learning and inference are one relation --- and it is
sharp only against the background this section describes.
\end{remark}

% ===================================================================
\section{What a coherent path would require}
\label{psy:sec:conditions}

The historical survey yields a specification.
Most of the conditions below name something an earlier tradition
supplied and its reception discarded; the last names something no
tradition in that survey supplied at all.

\begin{requirement}[Structural, not bit-level]
\label{psy:req:structural}
The representation the agent holds of itself must stand in a
disciplined relation to what it represents --- an explicit
quote/unquote pair with laws --- rather than being an unstructured
encoding.
Von Neumann's stored program failed here, and the failure is why the
practice became disreputable.
\end{requirement}

\begin{requirement}[Available at runtime, from inside]
\label{psy:req:runtime}
Reflection that requires an external tool-chain, an offline pass, or
a privileged operator is the observer's reflection under another
name.
The separation of training from inference
(Section~\ref{psy:sec:nn}) is the current and most consequential
violation of this condition, and it is architectural rather than
accidental.
Smith's criterion --- access to the environment and the
continuation, at the moment of processing --- is the right one.
\end{requirement}

\begin{requirement}[Priceable]
\label{psy:req:priced}
A regress cannot be budgeted; a loop can.
Self-modification without cost accounting is either unbounded or
arbitrarily curtailed, and neither is a mechanism.
Linear logic supplies the logical ancestor of this condition, but at
the level of the proof rather than of a running agent; what is
needed is resource sensitivity for the reflective act itself.
Note that pricing is not the property of any particular calculus:
the cost-accounting construction of Chapter~\ref{ch:carho} lifts to
an endofunctor on the category of interactive graph-structured
lambda theories, so this condition names a construction applicable
to a term language rather than a feature possessed by one.
Mortality is not a decoration on this list: an agent that can run
out is an agent for which reflective expenditure is a real decision.
\end{requirement}

\begin{requirement}[Social, and internally concurrent]
\label{psy:req:social}
The agent must be able to hold a representation of \emph{another}
agent's behavior, or reflection remains solitary and the resulting
system cannot participate in the peer relations that Petri, Hewitt
and Milner established as the real setting of computation.
And its own contact with the world must not be serialized: an agent
whose perception is a single queue has one sense organ whatever the
variety of its message types.
Rosette met the first requirement and, being built on actors, failed
the second.
\end{requirement}

\begin{requirement}[Self-supplying at the point of contact]
\label{psy:req:contact}
The means of contact must be constituted within the model rather
than left as an unanalyzed primitive to be filled in during
implementation.
A model whose ontology must be relaxed in order to be realized does
not determine what its agents perceive with; the implementer does.
And an agent cannot construct, modify or reason about an organ made
of material the calculus does not describe.
The general form of this condition is the ontological homogeneity of
Observation~\ref{psy:obs:isolation}: computation is isolated enough
that a model may be built from one sort, and every second sort a
model admits is a place where something must be supplied from
outside it.
\end{requirement}

\begin{requirement}[Closed rather than regressive]
\label{psy:req:closed}
Self-relation must terminate as a structure even where computation
does not, or every level of self-modification implies an unbuilt
level above it.
\end{requirement}

These conditions are jointly satisfiable, and --- more usefully ---
largely \emph{separable}.
Requirement~\ref{psy:req:priced} is a construction that can be
applied to a term language which does not yet have it.
Requirements~\ref{psy:req:structural}, \ref{psy:req:social}
and~\ref{psy:req:closed} are properties of the term language itself.
Requirements~\ref{psy:req:runtime} and~\ref{psy:req:contact} concern
the discipline relating a representation to what it represents.
A calculus in which names are quoted processes supplies five of the
six by construction and admits the sixth by the cost lifting, but
the point of stating them as a list is to make them separately
checkable, and separately obtainable, rather than to advertise a
solution.
The substantive prediction is that recursive self-improvement in any
strong sense will arrive, if it arrives, on a substrate meeting
Requirements~\ref{psy:req:structural}--\ref{psy:req:closed}, and not
by scaling a stack that meets none of them.

% ===================================================================
\section{Objections and limits}
\label{psy:sec:objections}

\textbf{Models reflect the question posed, not the psyche of the
poser.}
This is correct and is the strongest objection to the
individual-level version of the thesis.
Turing personally had a far richer conception than his machine
reveals: by 1948 he was writing about unorganized machines,
education, and the need for senses and a body \cite{turing1948}; by
1950 his proposed criterion for mind was not a computation but a
conversation, complete with deception and embarrassment
\cite{turing1950}.
The man who gave us the agent without an interlocutor made the
interlocutor the criterion of mind fourteen years later.
Von Neumann shipped the serial architecture while holding the
parallel, reproducing one.
The objection succeeds against individual psychologizing --- and in
succeeding, it establishes the community-level reading, since what
it demonstrates is precisely that the richer options were available
and were not taken up.
The Entscheidungsproblem received a formalism, an industry and a
hierarchy; the imitation game received a parlor argument.

\textbf{Practitioners can see the grounding and change it.}
Deutsch's regrounding of the Church--Turing thesis in physics rather
than in the capacities of a human clerk
(Section~\ref{psy:sec:quantum}) shows that the anthropomorphic basis
of a founding abstraction is not invisible to the people working
with it.
This is a real constraint on how much the method may claim:
communities are not simply enacting a psychology they cannot
perceive.
What the case also shows is that noticing the grounding does not by
itself enlarge the agent.
The replacement was chosen for physical fidelity, and it produced a
more solitary bot than the one it replaced.

\textbf{The transmission story is post hoc.}
Any historical sequence can be narrated as loss.
The check is that the selection criterion be stated in advance and
be falsifiable, which here it is: reception should favor the
operationally tractable half and lapse the half conferring standing
on the agent.
A counterexample would be a case where a field adopted the
agent-standing half at operational cost.
Candidates worth examining --- garbage collection, capability
security, session types, the partial rehabilitation of the
metaobject protocol in dynamic language runtimes --- are not
obviously counterexamples, but the analysis here is not deep enough
to rule them out, and someone should look.

\textbf{The reflection tradition's own self-criticism.}
Smith came to hold, by \emph{On the Origin of Objects}
\cite{smith1996}, that the computational tradition, his own work
included, had presupposed a world already carved into objects and
handed to the agent pre-individuated.
If that is right, then reflection was self-knowledge purchased
against an ontology the agent never negotiated, and the deepest
poverty in the survey is neither the missing other nor the missing
self but the pre-parcelled world.
This is the objection with the most force remaining, and it is the
one this book has the most direct answer to and the least complete
one.
The direct answer is Chapter~\ref{ch:sci}: an agent that pays for
its distinctions has at least begun to make them rather than inherit
them.
The incompleteness is that the distinctions it pays for are drawn
from a formula language it was given, which is
Section~\ref{psy:sec:openchallenges}'s first item and is not
answered anywhere in these pages.

% ===================================================================
\section{How life-like is enough?}
\label{psy:sec:threshold}

Stated positively rather than as a list of things receptions
dropped, the six conditions describe an agent that is progressively
more like a living thing.
It has peers and can represent them.
It is concurrent within itself, not merely with respect to its
surroundings.
Its organs of contact are made of its own material.
It is changed by what happens to it.
It can run out.
The thesis this chapter ends on is that this is the relevant axis.

\begin{principle}[Life-likeness and reach]
\label{psy:prin:lifelike}
The more life-like the phenomenology a model of computation
stipulates for its agents, the more available recursive
self-improvement becomes, and the further the model reaches toward
general intelligence.
\end{principle}

The argument is the one already made in
Section~\ref{psy:sec:present}.
The loop closes only when the improving process and the improved
process are the same process, in the same world, with the
improvement operation reachable from inside it; and the properties
that make that true are, one after another, the properties that make
an agent resemble an organism rather than an instrument.
Observation~\ref{psy:obs:three} is why: the two rules life breaks
are the two rules whose breaking makes improvement possible at all,
and priced discarding is the only one of the three positions at
which the criterion of improvement becomes internal to the agent.

\textbf{What is not known is where the boundary lies.}
``Life-like enough'' is not currently a quantity, and this chapter
cannot make it one.
Worse, the shape of the answer is unknown, not merely its value.
It is not established whether there is a threshold, a phase
transition, or a continuum with no boundary at all; whether the six
conditions are jointly required or whether some subset suffices; or
whether the list is complete.
Anyone claiming to know where the line falls is claiming more than
the evidence assembled here supports, and i am not exempt.

The contribution is therefore the decomposition and not the value.
The six conditions turn a gestalt into separable properties, each
independently obtainable and independently deniable, which is what
makes an ablation program possible at all: build systems possessing
$k$ of the six and find which removals break the closure of the
loop.
That is an empirical question.
Reading a century of receptions can generate such a question and
cannot settle it, and the gap between a historical instrument and an
empirical claim should be stated rather than finessed.

\begin{remark}[The book has one candidate quantity, and it is not
enough]
\label{psy:rem:assembly}
One candidate for a quantitative form of the question deserves
naming, and the reader met it three hundred pages ago.
Assembly theory measures how much construction history an object
requires rather than its instantaneous complexity
\cite{walker2023}, and Theorem~\ref{thm:ai-bound} bounds
discriminating power by the assembly index of the ecology that
supplies it.
So the book does turn a version of ``how much is enough'' into a
number carrying a unit.
Two things should be said about that before it is taken for more
than it is.
The number bounds \emph{resolution} --- how finely a population can
tell states apart --- and not reach, which
Remark~\ref{rmk:power-caveat} already says.
And nothing connects it to the six conditions of
Section~\ref{psy:sec:conditions}: a bound on what an ecology can
discriminate is not a bound on whether the improvement loop closes
inside it.
Whether assembly index is the right measure is open; that it
measures accumulated construction is the reason to look, since
accumulated construction is what the six conditions are all,
obliquely, about.
\end{remark}

\subsection{Open challenges for the symbolic line}
\label{psy:sec:openchallenges}

\begin{enumerate}[label=(\alph*),leftmargin=2.2em,itemsep=0.2em]
\item \textbf{Transduction.}
A reflective calculus has no account of where its categories come
from.
Proposing a trained network as the transducer is a division of
labor, not a solution, and it is Smith's objection
(Section~\ref{psy:sec:objections}) in different clothing: a world
handed over pre-individuated is a world whose symbols were handed
over too.
This is the same admission Chapter~\ref{ch:sci} makes in its own
vocabulary when it concedes that the network proposes and the
ecology only disposes.
\item \textbf{There is no gradient.}
Backpropagation works partly because credit assignment has a cheap
differentiable solution.
Search over a hypothesis space with an ultrametric structure, where
a single learner cannot leave its ball and escape requires
population, recombination and death, is biologically right and
computationally brutal.
The architecture may be correct and hopeless.
What plays the role of the gradient is, so far as i know,
unanswered.
\item \textbf{The environment is still stipulated.}
Death as failure to fund is endogenous only relative to an energy
economy somebody supplies.
If the rates are hand-tuned, the operator has returned wearing an
ecologist's hat.
Making the environment out of other agents pushes the stipulation
outward rather than removing it.
\item \textbf{Nothing has been demonstrated.}
The six conditions are proposed as necessary conditions with no
sufficiency argument, and no system meeting them has been shown to
close the loop.
The claim is a prediction, and its falsifier should be named: a
system meeting none of the conditions that recursively
self-improves anyway.
\item \textbf{The scale asymmetry.}
The connectionist line has an empirical scaling story.
This one has simulations.
Section~\ref{psy:sec:present} establishes that the current stack
cannot close the loop without people in it.
It does not establish that anything else can.
\end{enumerate}

\subsection{Four questions for the next reception}
\label{psy:sec:questions}

It would be inconsistent to spend a chapter arguing that reception
keeps the operational half and lapses the half conferring standing
on the agent, and then decline to turn the instrument on the work
this book is written from.

An earlier version of this material did that by making a forecast:
that what would travel from a cost-accounted reflective calculus is
the metering --- billing, resource limits, auditable execution ---
and that what would lapse is the reflection that made the metering
mean anything.
The concurrency would travel and the namespace discipline would not.
The weights would travel and the learner would not.

That forecast is an overreach and i have withdrawn it.
The instrument of Chapter~\ref{ch:psy} is built to read receptions
that have already happened; it establishes a pattern across cases
where the outcome is on the record and the counterfactuals are
gone.
Running it forwards on one's own work converts a historical
instrument into a prophecy, and does so in the one case where the
author has the least standing to be believed and the most reason to
be self-serving --- a forecast of one's own neglect is the cheapest
form of insurance a writer can buy.
What survives the withdrawal is better: the same content, stated as
questions a reader can go and check rather than as a prediction that
cannot be wrong in any way that costs anything.

\begin{question}
\label{psy:q:metering}
Does the metering travel without the reflection?
Cost accounting, resource limits and auditable execution are
separable from the calculus that motivated them, and
Requirement~\ref{psy:req:priced} says so explicitly.
If a decade from now the ledger is in wide use and the reflective
structure it was built to price is not, that is the pattern of
Principle~\ref{psy:prin:lapse} operating on this book.
If the two travel together, it is a counterexample, and the first
one this chapter has.
\end{question}

\begin{question}
\label{psy:q:concurrency}
Does concurrency travel without namespace discipline?
Parallel composition has already been received, several times over,
by industries that have no use for what a name is made of.
Chapter~\ref{ch:notation}'s account of notation is the part at risk
here, and the check is whether anything downstream of this book
treats the question of what a channel is made of as a question at
all.
\end{question}

\begin{question}
\label{psy:q:learner}
Do the weights travel without the learner?
The mortal scientist of Chapter~\ref{ch:sci} is a learner defined by
an account it maintains and can exhaust.
It is straightforward to extract from it a set of update rules and
discard the ledger, the mortality and the population, and the result
would still be publishable and would still work on the benchmarks
the field has.
Whether anyone does is the sharpest available test of the
selection criterion, because in this case the operational half and
the standing-conferring half are unusually easy to separate.
\end{question}

\begin{question}
\label{psy:q:counterexample}
Is there a reception that went the other way?
Section~\ref{psy:sec:objections} names four candidates and does not
resolve any of them.
A single well documented case in which a field paid an operational
price to keep the half that confers standing on the agent would not
refute Principle~\ref{psy:prin:lapse}, since it is a claim about a
tendency, but it would bound it, and nothing in this book bounds it
at all.
\end{question}

Naming the mechanism is the only known countermeasure, which is the
practical reason for writing a chapter of this kind at all.
Four questions with checkable answers are a better countermeasure
than one prediction with none.

% ===================================================================
\section{What is claimed}
\label{psy:sec:psy-claims}

Claimed: that a model of computation stipulates a phenomenology and
can be read for it; that reception, not invention, is the
instructive variable; that across a century the received half is the
operational one and the lapsed half is the one conferring standing
on the agent; that the contemporary stack is a metalevel constituted
entirely as an instrument of the operator; that in consequence there
is no recursive self-improvement in it without a person closing
every turn of the loop; and that the six conditions of
Section~\ref{psy:sec:conditions} name what a path would require.

Not claimed: that meeting the conditions is sufficient; that the
list is complete; that where the line falls is known; or that the
calculus this book is built on is the only thing that could meet
them.
Chapter~\ref{ch:psy} showed that the third generation was passed
over for operational reasons rather than refuted.
That is a fact about a community.
It is not, by itself, an argument that the third generation is
right, and the argument that it is right is the rest of this book.

There is one thing this pair of chapters does for the argument that
nothing else in the book could do, and it is worth naming before
Chapter~\ref{ch:kant} gathers everything.
Every other limit the book has found resolves into a price: what an
agent can distinguish, how deep it can look, what it can afford to
remember.
This one does not look like a price at first.
It looks like an absence --- a path that was never built.
But an absence with a hundred-year history of operational reasons
behind it is not a prohibition either.
It is the accumulated cost of a great many decisions each of which
was locally cheap, which is the same shape one level up, and it is
the last instance Chapter~\ref{ch:kant} will need.
